% Generated draft: requires mathematical and source review.
\documentclass[11pt]{article}
\usepackage[margin=25mm]{geometry}
\usepackage{fontspec}
\setmainfont{Latin Modern Roman}
\usepackage{amsmath,amssymb,mathtools,unicode-math}
\setmathfont{Latin Modern Math}
\usepackage{enumitem,needspace,xurl,hyperref}
\setlength{\parindent}{0pt}
\setlength{\parskip}{5pt}
\setlength{\emergencystretch}{3em}
\newcommand{\N}{\mathbb{N}}
\newcommand{\Z}{\mathbb{Z}}
\newcommand{\Q}{\mathbb{Q}}
\newcommand{\R}{\mathbb{R}}
\newcommand{\C}{\mathbb{C}}
\newcommand{\F}{\mathbb{F}}
\newcommand{\A}{\mathbb{A}}
\newcommand{\PP}{\mathbb{P}}
\newcommand{\E}{\mathbb{E}}
\newcommand{\eps}{\varepsilon}
\newcommand{\dd}{\,\mathrm{d}}
\begin{document}
\section*{UnsolvedMath ID: 30003677}
\textbf{Source problem code:} catalogue entry
\subsection{Definitions and mathematical statement}
For a finite vertex set \(V_n=\{1,\ldots,n\}\) and nonnegative integer degrees \(d_1^{(n)},\ldots,d_n^{(n)}\) with even sum, the configuration model pairs all half-edges uniformly at random to produce a random multigraph \(G_n\). Whether the intended model is conditioned to be simple is unspecified. The source assumes a finite-variance configuration model but does not specify a degree distribution, a deterministic degree-sequence regime, or connectivity assumptions. For a random degree variable \(D\), the usual finite-variance condition is \(\E[D^2]<\infty\); the precise asymptotic degree assumptions intended here remain unresolved.

In competing first-passage percolation, two types spread from their initial vertices along edges according to prescribed passage-time laws, and a vertex acquired by one type remains unavailable to the other. The passage-time laws, their dependence, and the rule for simultaneous arrivals must be specified to define the process. Write \(w\) for the type called weaker in the source and \(s\) for the type called stronger. These labels are retained without assigning rates or a stochastic ordering, since the source supplies neither. Let \(S_{w,n}\subseteq V_n\) be the initial vertices of type \(w\), and let \(v_{s,n}\in V_n\setminus S_{w,n}\) be the single initial vertex of type \(s\). Let \(C_{w,n}\) be the vertices eventually acquired by type \(w\), and put \(W_n=|C_{w,n}|\).

\textbf{Source question.} Can type \(w\) acquire a positive fraction of all vertices by starting at high-degree vertices, while type \(s\) starts at a low-degree vertex, in a finite-variance configuration model?

A possible asymptotic meaning of a positive fraction is
\[
\exists\,\varepsilon>0\quad
\liminf_{n\to\infty}\mathbb{P}(W_n\geq \varepsilon n)>0.
\]
This display is a provisional interpretation, not a recovered source quantifier. The source does not distinguish this requirement from success with probability tending to one or another probabilistic criterion. It also does not specify the number of weaker seeds, the degree thresholds, their dependence on \(n\), or how seeds are selected. Consequently, the supplied question does not yet determine a unique quantified mathematical proposition. Its intended target is whether a degree advantage at the initial placement can compensate for the weaker type's spreading disadvantage, once those missing parameters and hypotheses are fixed.
\subsection{Short English statement}
Can a slower or otherwise weaker infection occupy a positive fraction of a configuration-model graph by starting at high-degree vertices, when its stronger competitor starts at one low-degree vertex? The meanings of weaker, high degree, low degree, and probabilistic success are unspecified, so the formulation requires human mathematical review.
\subsection{Sources}
\begin{itemize}
\item Local UnsolvedMath registry, supplied complete record: permanent ID \(30003677\), source problem code \texttt{catalogue entry}, title \emph{Strategic Starting Vertices in Competing First-Passage Percolation}; collection \emph{Oberwolfach Reports Open Problems}. The supplied proposed year is 2017; no individual proposer is named.
\item Supplied report citation: \emph{Network Models: Structure and Function (2018)}, Oberwolfach Reports, DOI 10.4171/owr/2017/57, pp. 3427--3470. \url{https://doi.org/10.4171/owr/2017/57}. The background attributes the problem to workshop open-problem discussions documented in Oberwolfach Reports.
\item Supplied literature source: \emph{Competing first passage percolation on random graphs}. \url{https://arxiv.org/abs/1711.02902}.
\item Supplied literature source: \emph{OWR Report 57/2019}. \url{https://ems.press/content/serial-article-files/46721}.
\item Literature assessment dated exactly \texttt{2026-08-21}: the supplied curation labels the problem open and states that known competing first-passage results treat uniformly random starts and related regimes, while the high-degree-versus-low-degree initial placement question remains a stated further problem. This is supplied dated curation, not an independently verified literature assessment; its status and cited claims require independent verification.
\item The registry supplies the editorial label \texttt{corrected\_verified} and says that mathematical content was preserved while notation, typography, and sentence structure were repaired. These are supplied editorial assertions, not independent verification of this draft. The relationship between the report's supplied 2018 title date, its 2017 DOI component and proposed year, and the literature reference labelled 57/2019 remains unresolved.
\end{itemize}
\end{document}
